\chapter{系统测试与分析}
\label{chap:test}

\section{引言}
\label{sec:test-intro}

本章对匿踪纵向联邦学习系统进行全面的测试验证，包括功能测试、性能测试和模型效果评估。测试的目的是验证系统的正确性、安全性和性能，确保系统能够满足实际应用的需求。

测试工作分为以下几个部分：首先是测试环境配置，说明测试所使用的硬件、软件和网络环境；然后是单元测试，验证各模块内部逻辑的正确性；接着是功能测试，验证系统各功能模块是否按照需求规格正确运行；之后是性能测试，评估系统在不同数据规模和配置下的运行效率；最后是模型效果评估，在公开数据集上验证训练出的模型效果。

本章的测试用例设计对应第三章需求分析中定义的功能需求，确保每个核心功能都得到充分验证。

\section{测试准备}
\label{sec:test-prep}

\subsection{测试目的}
\label{subsec:test-purpose}

本次测试的主要目的包括以下几个方面：

\textbf{（1）功能正确性验证}：验证系统的各项功能是否按照需求规格正确实现，包括匿踪对齐、安全训练、模型推理等核心功能。功能测试确保系统能够正确完成端到端的纵向联邦学习任务。

\textbf{（2）性能评估}：评估系统在不同数据规模下的运行效率，包括PSI对齐时间、模型训练时间、推理延迟等关键性能指标。性能测试为系统的实际部署提供参考数据。

\textbf{（3）模型效果验证}：验证在隐私保护前提下，训练出的模型效果与非隐私保护方法相比是否具有可比性。模型效果测试评估隐私保护机制对模型精度的影响。

\subsection{测试环境}
\label{subsec:test-env}

测试环境的配置对测试结果的准确性和可重复性至关重要。本节详细说明测试所使用的硬件环境、软件环境、网络环境和测试工具。

\subsubsection{硬件环境}

测试使用的硬件环境配置如表~\ref{tab:test-hardware}所示。由于秘密共享计算需要在内存中维护多个数据分片，且密码学运算计算密集，测试环境配置了较大的内存和多核CPU。

\begin{table}[htbp]
    \centering
    \caption{测试硬件环境配置}
    \label{tab:test-hardware}
    \begin{tabular}{lp{8cm}}
        \toprule
        \textbf{配置项} & \textbf{详细说明} \\
        \midrule
        处理器 & Intel Core i9-12900，16核24线程，基础频率2.4GHz \\
        内存 & 32GB DDR4 3200MHz \\
        存储 & 512GB NVMe SSD \\
        网络 & 千兆以太网（局域网测试环境） \\
        操作系统 & Ubuntu 22.04 LTS \\
        \bottomrule
    \end{tabular}
\end{table}

在分布式测试场景下，使用了两台配置相同的服务器，一台部署Company方和Coordinator方、另一台部署Partner方。服务器之间通过局域网连接，网络延迟小于1毫秒。

\subsubsection{软件环境}

测试使用的软件环境如表~\ref{tab:test-software}所示。系统基于Python语言开发，采用SecretFlow框架作为隐私计算基础设施。

\begin{table}[htbp]
    \centering
    \caption{测试软件环境配置}
    \label{tab:test-software}
    \begin{tabular}{llp{5cm}}
        \toprule
        \textbf{类别} & \textbf{软件/库} & \textbf{版本及说明} \\
        \midrule
        开发语言 & Python & 3.10 \\
        \midrule
        \multirow{4}{*}{核心框架} & SecretFlow & 1.12.0b0，隐私计算框架 \\
        & Ray & 2.9.0，分布式计算框架 \\
        & JAX & 0.4.20，数值计算库 \\
        & NumPy & 1.26.0，数组操作库 \\
        \midrule
        密码学库 & rbcl & 1.0.1，Ristretto255椭圆曲线库 \\
        \midrule
        \multirow{2}{*}{测试工具} & pytest & 7.4.0，单元测试框架 \\
        & pytest-cov & 4.1.0，代码覆盖率插件 \\
        \midrule
        \multirow{2}{*}{数据处理} & Pandas & 2.1.0，数据读取和处理 \\
        & scikit-learn & 1.3.0，数据集划分和评估 \\
        \bottomrule
    \end{tabular}
\end{table}

pytest是Python语言广泛使用的单元测试框架，支持测试发现、夹具（fixture）、参数化等特性。pytest-cov插件用于统计测试的代码覆盖率，帮助评估测试的完整性。

\subsubsection{网络环境}

测试的网络环境配置如下：

\textbf{单机仿真模式}：所有参与方在同一台机器的同一Python进程中运行，通过内部通信机制进行数据交换，不涉及实际网络传输。

\textbf{局域网多机模式}：三台服务器部署在同一局域网内，通过千兆交换机连接。各节点的网络地址配置如表~\ref{tab:test-network}所示。

\begin{table}[htbp]
    \centering
    \caption{局域网测试网络配置}
    \label{tab:test-network}
    \begin{tabular}{lll}
        \toprule
        \textbf{节点角色} & \textbf{IP地址} & \textbf{SPU端口} \\
        \midrule
        Company & 192.168.1.100 & 9394 \\
        Partner & 192.168.1.101 & 9395 \\
        Coordinator & 192.168.1.100 & 9396 \\
        \bottomrule
    \end{tabular}
\end{table}

\subsubsection{测试工具}

测试过程中使用的主要工具包括：

\textbf{pytest}：用于编写和执行单元测试，支持自动发现测试用例、生成测试报告。

\textbf{pytest-cov}：用于统计代码覆盖率，生成覆盖率报告。

\textbf{time模块}：用于测量各阶段运行时间，评估性能。

\textbf{memory\_profiler}：用于监控内存使用情况，分析内存开销。

\textbf{TensorBoard}（可选）：用于可视化训练过程中的指标变化。

\section{单元测试}
\label{sec:unit-test}

单元测试是软件测试的基础环节，用于验证各模块内部逻辑的正确性。本系统的单元测试采用pytest框架，对公共模块、PSI模块、LR模块和XGBoost模块的核心功能进行了测试。

\subsection{测试框架说明}

系统使用pytest作为单元测试框架。pytest具有以下优点：语法简洁，使用assert语句进行断言；支持自动发现测试用例（以test\_开头或\_test结尾的文件和函数）；支持夹具（fixture）机制，方便测试环境的搭建和清理；支持参数化测试，可以用同一测试逻辑测试多组输入；丰富的插件生态，支持覆盖率统计、并行执行等功能。

测试代码组织在tests目录下，与源代码目录结构对应。每个模块的测试代码放在独立的测试文件中，便于管理和维护。

\subsection{单元测试结果}

各模块的单元测试结果如表~\ref{tab:unit-test-result}所示。测试覆盖了各模块的核心功能，包括激活函数计算、损失函数计算、梯度计算、模型前向传播和反向传播等。

\begin{table}[htbp]
    \centering
    \caption{单元测试结果统计}
    \label{tab:unit-test-result}
    \begin{tabular}{lcccc}
        \toprule
        \textbf{模块名称} & \textbf{测试用例数} & \textbf{通过数} & \textbf{通过率} & \textbf{代码覆盖率} \\
        \midrule
        公共模块（common） & XX个 & XX个 & XX\% & XX\% \\
        PSI模块 & XX个 & XX个 & XX\% & XX\% \\
        LR模块 & XX个 & XX个 & XX\% & XX\% \\
        XGBoost模块 & XX个 & XX个 & XX\% & XX\% \\
        推理模块 & XX个 & XX个 & XX\% & XX\% \\
        \midrule
        \textbf{合计} & \textbf{XX个} & \textbf{XX个} & \textbf{XX\%} & \textbf{XX\%} \\
        \bottomrule
    \end{tabular}
\end{table}

\subsection{公共模块测试}

公共模块的单元测试主要验证不同部署模式下的计算安全多方计算环境初始化正确性，包括SPU设备，company、partner、coordinator三个PYU设备以及两套HEU设备的正确创建。

\subsection{PSI模块测试}

PSI模块的单元测试验证PS\textsuperscript{3}I协议的正确性。

\textbf{哈希函数测试}：验证SHA512哈希到椭圆曲线点的映射是确定性的，相同的输入产生相同的输出。

\textbf{标量乘法测试}：验证椭圆曲线标量乘法的正确性，确保$k \cdot P$的计算结果符合预期。

\textbf{密钥交换测试}：验证Diffie-Hellman密钥交换的正确性，确保$k_c \cdot k_p \cdot P = k_p \cdot k_c \cdot P$。

\textbf{交集计算测试}：使用已知的测试数据验证交集计算的正确性，确保双方共同用户的哈希值能够正确匹配。

\subsection{LR模块测试}

LR模块的单元测试验证秘密共享逻辑回归算法的正确性。

\textbf{前向传播测试}：验证线性变换$z = X \cdot w$的计算结果与NumPy计算结果一致；验证激活函数应用后的输出范围正确。

\textbf{反向传播测试}：验证梯度计算公式$\frac{\partial L}{\partial w} = X^T \cdot (\hat{y} - y)$的正确性；验证权重更新公式的正确性。

\textbf{模型持久化测试}：验证模型保存和加载的正确性，确保加载后的模型能够产生与原模型相同的预测结果。

\subsection{XGBoost模块测试}

XGBoost模块的单元测试验证秘密共享XGBoost算法的正确性。

\textbf{分桶量化测试}：验证等频分桶的正确性，确保每个桶内的样本数量大致相等；验证桶标签的正确性。

\textbf{梯度聚合测试}：验证桶内梯度求和的正确性，确保聚合结果与在明文上逐元素求和结果一致。

\textbf{增益计算测试}：验证增益公式的正确性；验证无除法比较方法的正确性，确保与直接计算增益比较的结果一致。

\textbf{决策树构建测试}：验证单棵决策树的构建过程，确保分裂点选择和叶子节点权重计算正确。

\section{功能测试}
\label{sec:func-test}

功能测试验证系统的各项功能是否按照需求规格正确运行。本节按照匿踪对齐、安全训练、模型推理三个核心功能进行测试。

\subsection{测试用例设计}

功能测试用例的设计对应第三章需求分析中定义的功能需求。表~\ref{tab:func-test-cases}列出了功能测试用例的基本信息。

\begin{table}[htbp]
    \centering
    \caption{功能测试用例基本信息}
    \label{tab:func-test-cases}
    \begin{tabular}{llll}
        \toprule
        \textbf{用例ID} & \textbf{对应需求} & \textbf{测试名称} & \textbf{测试结果} \\
        \midrule
        TC1 & FR-01 & 匿踪对齐功能测试 & 待测试 \\
        TC2 & FR-03 & SSLR模型训练测试 & 待测试 \\
        TC3 & FR-04 & SSXGBoost模型训练测试 & 待测试 \\
        TC4 & FR-05 & 模型保存与加载测试 & 待测试 \\
        TC5 & FR-06 & 模型推理功能测试 & 待测试 \\
        \bottomrule
    \end{tabular}
\end{table}

\subsection{匿踪对齐功能测试}

匿踪对齐是系统的核心功能，测试用例TC1验证PS\textsuperscript{3}I协议的正确性。

\subsubsection{测试用例TC1：匿踪对齐功能测试}

测试用例TC1的详细信息如表~\ref{tab:tc1-psi}所示。

\begin{table}[htbp]
    \centering
    \caption{测试用例TC1：匿踪对齐功能测试}
    \label{tab:tc1-psi}
    \begin{tabular}{lp{9cm}}
        \toprule
        \textbf{项目} & \textbf{内容描述} \\
        \midrule
        测试ID & TC1 \\
        测试名称 & 匿踪对齐功能测试 \\
        测试功能 & 验证PS\textsuperscript{3}I协议能够正确完成样本对齐，且不泄露非交集成员信息 \\
        \midrule
        \multirow{5}{*}{测试步骤} & 1. 准备Company方和Partner方的测试数据，包含已知的交集用户和非交集用户； \\
        & 2. 调用PSI模块执行匿踪对齐； \\
        & 3. 验证输出分片的样本数量等于交集大小； \\
        & 4. 验证输出分片的特征值正确（分片之和等于原始特征值）； \\
        & 5. 验证Company方无法获取Partner方的非交集用户信息 \\
        \midrule
        预期结果 & 对齐后的样本数量正确，特征分片正确，非交集成员信息不泄露 \\
        实际结果 & 待测试 \\
        \bottomrule
    \end{tabular}
\end{table}

\subsection{安全训练功能测试}

安全训练功能测试验证SSLR和SSXGBoost两种模型的训练正确性。

\subsubsection{测试用例TC2：SSLR模型训练测试}

测试用例TC2的详细信息如表~\ref{tab:tc2-sslr}所示。

\begin{table}[htbp]
    \centering
    \caption{测试用例TC2：SSLR模型训练测试}
    \label{tab:tc2-sslr}
    \begin{tabular}{lp{9cm}}
        \toprule
        \textbf{项目} & \textbf{内容描述} \\
        \midrule
        测试ID & TC2 \\
        测试名称 & SSLR模型训练测试 \\
        测试功能 & 验证秘密共享逻辑回归模型能够正确训练并收敛 \\
        \midrule
        \multirow{6}{*}{测试步骤} & 1. 准备已对齐的秘密共享训练数据； \\
        & 2. 创建SSLR模型实例，配置训练参数； \\
        & 3. 调用fit方法执行训练； \\
        & 4. 监控训练过程中的损失变化和验证指标； \\
        & 5. 验证训练完成后模型权重的有效性； \\
        & 6. 使用训练好的模型进行预测，验证预测结果 \\
        \midrule
        预期结果 & 模型能够正常训练，损失逐渐下降，验证指标达到预期水平 \\
        实际结果 & 待测试 \\
        \bottomrule
    \end{tabular}
\end{table}

\subsubsection{测试用例TC3：SSXGBoost模型训练测试}

测试用例TC3的详细信息如表~\ref{tab:tc3-ssxgboost}所示。

\begin{table}[htbp]
    \centering
    \caption{测试用例TC3：SSXGBoost模型训练测试}
    \label{tab:tc3-ssxgboost}
    \begin{tabular}{lp{9cm}}
        \toprule
        \textbf{项目} & \textbf{内容描述} \\
        \midrule
        测试ID & TC3 \\
        测试名称 & SSXGBoost模型训练测试 \\
        测试功能 & 验证秘密共享XGBoost模型能够正确训练并收敛 \\
        \midrule
        \multirow{6}{*}{测试步骤} & 1. 准备已对齐的秘密共享训练数据和分桶信息； \\
        & 2. 创建SSXGBoost模型实例，配置训练参数； \\
        & 3. 调用fit方法执行训练； \\
        & 4. 监控每棵树训练后的验证指标； \\
        & 5. 验证训练完成后决策树结构的有效性； \\
        & 6. 使用训练好的模型进行预测，验证预测结果 \\
        \midrule
        预期结果 & 模型能够正常训练，验证指标随树数量增加而提升 \\
        实际结果 & 待测试 \\
        \bottomrule
    \end{tabular}
\end{table}

\subsubsection{测试用例TC4：模型保存与加载测试}

测试用例TC4验证模型持久化功能的正确性，详细信息如表~\ref{tab:tc4-persistence}所示。

\begin{table}[htbp]
    \centering
    \caption{测试用例TC4：模型保存与加载测试}
    \label{tab:tc4-persistence}
    \begin{tabular}{lp{9cm}}
        \toprule
        \textbf{项目} & \textbf{内容描述} \\
        \midrule
        测试ID & TC4 \\
        测试名称 & 模型保存与加载测试 \\
        测试功能 & 验证模型能够正确保存和加载，加载后的模型功能正常 \\
        \midrule
        \multirow{5}{*}{测试步骤} & 1. 训练SSLR和SSXGBoost模型； \\
        & 2. 使用训练好的模型对测试数据进行预测，记录预测结果； \\
        & 3. 调用save方法保存模型到指定路径； \\
        & 4. 调用load方法加载模型； \\
        & 5. 使用加载后的模型对相同的测试数据进行预测； \\
        & 6. 比较加载前后模型的预测结果是否一致 \\
        \midrule
        预期结果 & 加载后的模型预测结果与原模型完全一致 \\
        实际结果 & 待测试 \\
        \bottomrule
    \end{tabular}
\end{table}

\subsection{模型推理功能测试}

模型推理功能测试验证推理引擎的正确性。

\subsubsection{测试用例TC5：模型推理功能测试}

测试用例TC5的详细信息如表~\ref{tab:tc5-infer}所示。

\begin{table}[htbp]
    \centering
    \caption{测试用例TC5：模型推理功能测试}
    \label{tab:tc5-infer}
    \begin{tabular}{lp{9cm}}
        \toprule
        \textbf{项目} & \textbf{内容描述} \\
        \midrule
        测试ID & TC5 \\
        测试名称 & 模型推理功能测试 \\
        测试功能 & 验证推理引擎能够正确加载模型并执行预测 \\
        \midrule
        \multirow{5}{*}{测试步骤} & 1. 准备已训练的模型文件； \\
        & 2. 准备推理数据（Company方和Partner方各自的CSV文件）； \\
        & 3. 创建InferEngine实例，加载模型； \\
        & 4. 调用compute\_intersection方法对推理数据求交集并对齐； \\
        & 5. 调用infer方法执行推理； \\
        & 6. 验证推理输出的格式正确（包含ID和预测结果）； \\
        & 7. 验证推理结果的合理性 \\
        \midrule
        预期结果 & 推理输出格式正确，预测结果在合理范围内 \\
        实际结果 & 待测试 \\
        \bottomrule
    \end{tabular}
\end{table}

\section{性能测试}
\label{sec:perf-test}

性能测试评估系统在不同数据规模和配置下的运行效率。本节分别测试PSI性能、训练性能和推理性能。

性能测试采用以下方案：准备不同规模的随机测试数据集，样本数量从1000到100,000不等；特征维度固定为20维（Company方10维，Partner方10维）；分别测量各阶段的运行时间，记录通信数据量。

\subsection{PSI性能测试}

PSI性能测试评估匿踪对齐协议在不同样本规模下的运行时间和通信开销。

\subsubsection{测试结果}

PSI性能测试结果如表~\ref{tab:psi-perf}所示。

\begin{table}[htbp]
    \centering
    \caption{PSI性能测试结果}
    \label{tab:psi-perf}
    \begin{tabular}{cccc}
        \toprule
        \textbf{样本数量} & \textbf{对齐时间} & \textbf{通信数据量} & \textbf{内存占用} \\
        \midrule
        1,000 & 待测试 & 待测试 & 待测试 \\
        5,000 & 待测试 & 待测试 & 待测试 \\
        10,000 & 待测试 & 待测试 & 待测试 \\
        50,000 & 待测试 & 待测试 & 待测试 \\
        100,000 & 待测试 & 待测试 & 待测试 \\
        \bottomrule
    \end{tabular}
\end{table}

\subsubsection{结果分析}

% [待补充：PSI性能测试结果分析]
（待测试完成后补充）

\subsection{训练性能测试}

训练性能测试评估SSLR和SSXGBoost模型在不同配置下的训练时间。

\subsubsection{SSLR训练性能测试}

SSLR训练性能测试结果如表~\ref{tab:sslr-perf}所示。测试配置为：20个训练轮次，批次大小1024，学习率0.1。

\begin{table}[htbp]
    \centering
    \caption{SSLR训练性能测试结果}
    \label{tab:sslr-perf}
    \begin{tabular}{cccc}
        \toprule
        \textbf{样本数量} & \textbf{特征维度} & \textbf{单轮训练时间} & \textbf{总训练时间} \\
        \midrule
        1,000 & 20 & 待测试 & 待测试 \\
        5,000 & 20 & 待测试 & 待测试 \\
        10,000 & 20 & 待测试 & 待测试 \\
        50,000 & 20 & 待测试 & 待测试 \\
        100,000 & 20 & 待测试 & 待测试 \\
        \bottomrule
    \end{tabular}
\end{table}

\subsubsection{SSXGBoost训练性能测试}

SSXGBoost训练性能测试结果如表~\ref{tab:ssxgboost-perf}所示。测试配置为：10棵决策树，最大深度3，分位点数10。

\begin{table}[htbp]
    \centering
    \caption{SSXGBoost训练性能测试结果}
    \label{tab:ssxgboost-perf}
    \begin{tabular}{cccc}
        \toprule
        \textbf{样本数量} & \textbf{特征维度} & \textbf{单棵树训练时间} & \textbf{总训练时间} \\
        \midrule
        1,000 & 20 & 待测试 & 待测试 \\
        5,000 & 20 & 待测试 & 待测试 \\
        10,000 & 20 & 待测试 & 待测试 \\
        50,000 & 20 & 待测试 & 待测试 \\
        100,000 & 20 & 待测试 & 待测试 \\
        \bottomrule
    \end{tabular}
\end{table}

\subsubsection{批次大小对性能的影响}

批次大小对SSLR训练性能的影响如表~\ref{tab:batch-size-perf}所示。测试配置为：10万样本，20维特征，10个训练轮次。

\begin{table}[htbp]
    \centering
    \caption{批次大小对SSLR训练性能的影响}
    \label{tab:batch-size-perf}
    \begin{tabular}{ccc}
        \toprule
        \textbf{批次大小} & \textbf{单轮迭代次数} & \textbf{总训练时间} \\
        \midrule
        256 & 待测试 & 待测试 \\
        512 & 待测试 & 待测试 \\
        1024 & 待测试 & 待测试 \\
        2048 & 待测试 & 待测试 \\
        4096 & 待测试 & 待测试 \\
        \bottomrule
    \end{tabular}
\end{table}

\subsection{推理性能测试}

推理性能测试评估模型推理的延迟和吞吐量。

\subsubsection{推理延迟测试}

推理延迟测试结果如表~\ref{tab:infer-latency}所示。测试配置为：单次推理请求的样本数量从1到10000不等。

\begin{table}[htbp]
    \centering
    \caption{推理延迟测试结果}
    \label{tab:infer-latency}
    \begin{tabular}{ccc}
        \toprule
        \textbf{样本数量} & \textbf{推理延迟} & \textbf{单样本平均延迟} \\
        \midrule
        1 & 待测试 & 待测试 \\
        10 & 待测试 & 待测试 \\
        100 & 待测试 & 待测试 \\
        1,000 & 待测试 & 待测试 \\
        10,000 & 待测试 & 待测试 \\
        \bottomrule
    \end{tabular}
\end{table}

\subsubsection{吞吐量测试}

推理吞吐量测试结果如表~\ref{tab:infer-throughput}所示。

\begin{table}[htbp]
    \centering
    \caption{推理吞吐量测试结果}
    \label{tab:infer-throughput}
    \begin{tabular}{cc}
        \toprule
        \textbf{模型类型} & \textbf{吞吐量（样本/秒）} \\
        \midrule
        SSLR & 待测试 \\
        SSXGBoost（10棵树） & 待测试 \\
        SSXGBoost（50棵树） & 待测试 \\
        \bottomrule
    \end{tabular}
\end{table}

\subsection{性能测试结果分析}

% [待补充：性能测试结果分析]
（待测试完成后补充）

\section{模型效果评估}
\label{sec:model-eval}

模型效果评估验证在隐私保护前提下，训练出的模型效果与非隐私保护方法相比是否具有可比性。本节在公开数据集上进行实验，对比分析不同方法的模型效果。

\subsection{实验数据集}

\subsubsection{数据集描述}

实验使用以下公开数据集进行模型效果评估：

\textbf{（1）UCI Adult数据集}：该数据集包含美国人口普查数据，用于预测个人年收入是否超过50K美元。数据集包含48842个样本，14个特征（包括年龄、教育程度、职业等）。在纵向联邦学习场景下，将特征划分为两部分，分别由Company方和Partner方持有。

\textbf{（2）Give Credit数据集}：该数据集包含信贷申请数据，用于预测客户是否会违约。数据集包含（待补充）个样本，（待补充）个特征。

数据集的详细信息如表~\ref{tab:dataset-info}所示。

\begin{table}[htbp]
    \centering
    \caption{实验数据集信息}
    \label{tab:dataset-info}
    \begin{tabular}{lccc}
        \toprule
        \textbf{数据集} & \textbf{样本数量} & \textbf{特征数量} & \textbf{任务类型} \\
        \midrule
        UCI Adult & 48,842 & 14 & 二分类 \\
        Give Credit & 待补充 & 待补充 & 二分类 \\
        \bottomrule
    \end{tabular}
\end{table}

\subsubsection{数据预处理}

数据预处理包括以下步骤：

\textbf{（1）缺失值处理}：对于数值型特征，使用均值填充缺失值；对于类别型特征，使用众数填充缺失值。

\textbf{（2）特征编码}：类别型特征使用One-Hot编码转换为数值型特征。

\textbf{（3）特征标准化}：数值型特征进行标准化处理，使其均值为0，标准差为1。

\textbf{（4）数据集划分}：将数据集按8:1:1的比例划分为训练集、验证集和测试集。

\textbf{（5）特征划分}：将特征按纵向划分为两部分，模拟两个参与方各自持有部分特征的场景。

\subsubsection{用户ID模拟}

由于公开数据集不包含用户ID，实验中使用行索引模拟用户ID。为模拟真实场景中两方用户集合不完全重叠的情况，实验中对两方的数据进行采样，使得两方的用户集合存在部分重叠。

\subsection{评估指标}

模型效果评估使用以下指标：

\textbf{（1）准确率（Accuracy）}：正确预测的样本数占总样本数的比例。

\textbf{（2）精确率（Precision）}：预测为正例中实际为正例的比例。

\textbf{（3）召回率（Recall）}：实际为正例中被正确预测为正例的比例。

\textbf{（4）F1分数}：精确率和召回率的调和平均数。

\textbf{（5）AUC}：ROC曲线下的面积，衡量模型的整体分类能力。

\subsection{对比实验}

\subsubsection{基线方法}

对比实验包括以下基线方法：

\textbf{（1）明文训练（Plaintext）}：在明文状态下使用标准逻辑回归和XGBoost算法训练模型，不进行任何隐私保护。该方法代表理想情况下的模型效果上限。

\textbf{（2）近似SSLR（Approx SSLR）}：使用分段线性函数近似Sigmoid激活函数的秘密共享逻辑回归。

\textbf{（3）精确SSLR（Exact SSLR）}：使用标准Sigmoid函数，标签保持明文的秘密共享逻辑回归。

\textbf{（4）SSXGBoost}：秘密共享XGBoost算法。

\subsubsection{实验配置}

实验配置如表~\ref{tab:exp-config}所示。

\begin{table}[htbp]
    \centering
    \caption{实验配置}
    \label{tab:exp-config}
    \begin{tabular}{lll}
        \toprule
        \textbf{参数} & \textbf{SSLR配置} & \textbf{SSXGBoost配置} \\
        \midrule
        训练轮次/树数量 & 10 & 10 \\
        批次大小 & 1024 & - \\
        学习率 & 0.1 & 0.1 \\
        L2正则化系数 & 0.01 & 0.1 \\
        最大深度 & - & 3 \\
        分位点数 & - & 10 \\
        \bottomrule
    \end{tabular}
\end{table}

\subsubsection{实验结果}

在UCI Adult数据集上的实验结果如表~\ref{tab:exp-adult}所示。

\begin{table}[htbp]
    \centering
    \caption{UCI Adult数据集实验结果}
    \label{tab:exp-adult}
    \begin{tabular}{lccccc}
        \toprule
        \textbf{方法} & \textbf{准确率} & \textbf{精确率} & \textbf{召回率} & \textbf{F1分数} & \textbf{AUC} \\
        \midrule
        明文逻辑回归 & 待测试 & 待测试 & 待测试 & 待测试 & 待测试 \\
        近似SSLR & 待测试 & 待测试 & 待测试 & 待测试 & 待测试 \\
        精确SSLR & 待测试 & 待测试 & 待测试 & 待测试 & 待测试 \\
        \midrule
        明文XGBoost & 待测试 & 待测试 & 待测试 & 待测试 & 待测试 \\
        SSXGBoost & 待测试 & 待测试 & 待测试 & 待测试 & 待测试 \\
        \bottomrule
    \end{tabular}
\end{table}

在Give Credit数据集上的实验结果如表~\ref{tab:exp-credit}所示。

\begin{table}[htbp]
    \centering
    \caption{Give Credit数据集实验结果}
    \label{tab:exp-credit}
    \begin{tabular}{lccccc}
        \toprule
        \textbf{方法} & \textbf{准确率} & \textbf{精确率} & \textbf{召回率} & \textbf{F1分数} & \textbf{AUC} \\
        \midrule
        明文逻辑回归 & 待测试 & 待测试 & 待测试 & 待测试 & 待测试 \\
        近似SSLR & 待测试 & 待测试 & 待测试 & 待测试 & 待测试 \\
        精确SSLR & 待测试 & 待测试 & 待测试 & 待测试 & 待测试 \\
        \midrule
        明文XGBoost & 待测试 & 待测试 & 待测试 & 待测试 & 待测试 \\
        SSXGBoost & 待测试 & 待测试 & 待测试 & 待测试 & 待测试 \\
        \bottomrule
    \end{tabular}
\end{table}

\subsection{实验结果分析}

% [待补充：实验结果分析]
（待测试完成后补充）

\subsubsection{隐私保护与模型效果的权衡}

% [待补充：分析隐私保护机制对模型效果的影响]
（待测试完成后补充）

\subsubsection{近似模式与精确模式的对比}

% [待补充：分析近似SSLR和精确SSLR的效果差异]
（待测试完成后补充）

\section{本章小结}
\label{sec:test-summary}

本章对匿踪纵向联邦学习系统进行了全面的测试验证，包括单元测试、功能测试、性能测试和模型效果评估。

在测试准备阶段，详细说明了测试目的和测试环境配置，包括硬件环境、软件环境、网络环境和测试工具。

在单元测试阶段，使用pytest框架对各模块的核心功能进行了测试，验证了激活函数、损失函数、梯度计算、模型前向传播和反向传播等功能的正确性。

在功能测试阶段，设计了8个测试用例，覆盖了匿踪对齐、数据分片生成、安全训练、模型持久化和模型推理等核心功能。测试用例的设计对应第三章需求分析中定义的功能需求。

在性能测试阶段，测试了PSI性能、训练性能和推理性能。测试结果表明（待测试完成后补充）。

在模型效果评估阶段，在UCI Adult和Give Credit数据集上进行了实验，对比了明文训练、近似SSLR、精确SSLR和SSXGBoost等方法的效果。实验结果表明（待测试完成后补充）。

通过本章的测试工作，验证了匿踪纵向联邦学习系统的正确性、安全性和性能，为系统的实际应用提供了保障。